\documentclass[11pt,a4paper]{article}
\usepackage[utf8]{inputenc}
\usepackage[T1]{fontenc}
\usepackage[catalan]{babel}
\usepackage{lmodern}
\usepackage[sfdefault,lining]{FiraSans}
\usepackage[a4paper,top=2.3cm,bottom=2.2cm,left=1.8cm,right=1.8cm,headheight=1.3cm,headsep=0.4cm,footskip=1cm]{geometry}
\usepackage{amsmath,amssymb}
\usepackage{icomma}
\usepackage{graphicx}
\usepackage[table]{xcolor}
\usepackage{tikz}
\usepackage{fontawesome5}
\usepackage{ccicons}
\usepackage{fancyhdr}
\usepackage{lastpage}
\usepackage{tasks}
\usepackage{needspace}
\usepackage[most]{tcolorbox}
\usepackage{array}
\usepackage[hidelinks,pdftitle={Fitxa d'exercicis de la unitat 7: Funcions. 3r d'ESO},pdfsubject={Matemàtiques 3r ESO},pdfkeywords={funcions, imatge, antiimatge, taula de valors, gràfiques, domini, recorregut, creixement, màxims i mínims, continuïtat}]{hyperref}

\definecolor{navy}{HTML}{1B2748}
\definecolor{accent}{HTML}{E8603A}
\definecolor{teal}{HTML}{1F9E90}
\definecolor{soft}{HTML}{EEF1F8}
\definecolor{muted}{HTML}{5D6685}

\setlength{\parindent}{0pt}
\setlength{\parskip}{3pt}
\settasks{label=\alph*), label-format=\bfseries\color{navy}, label-width=1.3em, item-indent=1.8em, column-sep=0.6em, after-item-skip=5pt, before-skip=3pt, after-skip=4pt}

% ---------- capçalera i peu ----------
\newcommand{\llicencia}{\href{https://creativecommons.org/licenses/by-nc-sa/4.0/deed.ca}{\ccbyncsa}\ \ Llicència Creative Commons Reconeixement-NoComercial-CompartirIgual 4.0 Internacional (CC BY-NC-SA 4.0)}
\pagestyle{fancy}
\fancyhf{}
\fancyhead[L]{\raisebox{-0.35\height}{\includegraphics[height=0.95cm]{logo-gva.png}}}
\fancyhead[C]{\small\color{muted}Fitxa d'exercicis · Unitat 7: Funcions · 3r d'ESO}
\fancyhead[R]{\raisebox{-0.35\height}{\includegraphics[height=0.95cm]{logo-vermellar.png}}}
\fancyfoot[L]{\scriptsize\color{muted}\llicencia}
\fancyfoot[R]{\small\color{muted}\thepage}
\renewcommand{\headrulewidth}{0.4pt}
\renewcommand{\headrule}{\color{muted!40}\hrule height \headrulewidth}
\fancypagestyle{primera}{\fancyhf{}\renewcommand{\headrulewidth}{0pt}
  \fancyfoot[L]{\scriptsize\color{muted}\llicencia}
  \fancyfoot[R]{\small\color{muted}\thepage}}

% ---------- nivells i exercicis ----------
\newcommand{\estrelles}[1]{{\color{accent}\ifcase#1\or\faStar\faStar[regular]\faStar[regular]\or\faStar\faStar\faStar[regular]\or\faStar\faStar\faStar\fi}}
\newcounter{exer}
\newcommand{\ex}[1]{\par\needspace{4\baselineskip}\medskip\refstepcounter{exer}\noindent{\color{navy}\bfseries\theexer.}\ \ {\small\estrelles{#1}}\ \ }

\newcommand{\seccio}[2]{\par\needspace{10\baselineskip}\bigskip
  \begin{tcolorbox}[colback=navy,colframe=navy,coltext=white,arc=2mm,boxrule=0pt,left=3mm,right=3mm,top=1.2mm,bottom=1.2mm]
  {\large\bfseries #1\quad #2}
  \end{tcolorbox}\vspace{-2pt}}

\newcolumntype{C}[1]{>{\centering\arraybackslash}p{#1}}
\usepackage{pifont}

% quadre de recordatori (només a la fitxa d'exercicis)
\newcommand{\recorda}[1]{\par\smallskip\begin{tcolorbox}[colback=accent!7,colframe=accent!55,arc=2mm,boxrule=0.5pt,left=2.5mm,right=2.5mm,top=1mm,bottom=1mm]\small{\color{accent}\faLightbulb[regular]}\ \textbf{Recorda:}\ #1\end{tcolorbox}}
\newcommand{\res}[1]{#1}
\newcommand{\rest}[1]{#1}

\def\fmneed{14}
% =====================================================================
% Capçalera FentMates de les fitxes: un color per a cada curs.
% Els colors són els de la pàgina inicial (LEVELS a assets/js/materials.js);
% el logo ve de _eines/logo/pdf/ (python3 _eines/logo/genera.py).
%
% Ús, al final del preàmbul de la unitat:
%   % =====================================================================
% Capçalera FentMates de les fitxes: un color per a cada curs.
% Els colors són els de la pàgina inicial (LEVELS a assets/js/materials.js);
% el logo ve de _eines/logo/pdf/ (python3 _eines/logo/genera.py).
%
% Ús, al final del preàmbul de la unitat:
%   % =====================================================================
% Capçalera FentMates de les fitxes: un color per a cada curs.
% Els colors són els de la pàgina inicial (LEVELS a assets/js/materials.js);
% el logo ve de _eines/logo/pdf/ (python3 _eines/logo/genera.py).
%
% Ús, al final del preàmbul de la unitat:
%   \input{../comu/fentmates.tex}
%   \fmcurs{1eso}
% i, al començament del document:
%   \capcalera{Matemàtiques}{1r ESO}{Fitxa d'exercicis · Unitat 2}{Divisibilitat}
% Opcional abans de \input: \def\fmneed{14} (espai mínim abans de cada secció).
% =====================================================================
\makeatletter
\def\fm@logodir{../../logo/pdf}
\providecommand{\fmneed}{10}
\def\fm@cursnom{}

% clau del curs, color i color del text damunt del color
\newcommand{\fm@curs}[3]{\@namedef{fm@color@#1}{#2}\@namedef{fm@text@#1}{#3}}
\fm@curs{1eso}{E6B84A}{navy}
\fm@curs{2eso}{3B82F6}{white}
\fm@curs{3eso}{FF7A50}{white}
\fm@curs{4eso-a}{A78BFA}{white}
\fm@curs{4eso-b}{6D28D9}{white}
\fm@curs{1batx-mi}{9B1C1C}{white}
\fm@curs{1batx-ccss}{EF6B6B}{white}
\fm@curs{2batx-mii}{1E6B47}{white}
\fm@curs{2batx-ccss}{2BB3A3}{white}

\newcommand{\fmcurs}[1]{%
  \@ifundefined{fm@color@#1}{\PackageError{fentmates}{Curs desconegut: #1}{}}{}%
  \edef\fm@tmp{\noexpand\definecolor{curs}{HTML}{\@nameuse{fm@color@#1}}%
               \noexpand\colorlet{curstext}{\@nameuse{fm@text@#1}}}\fm@tmp
  \colorlet{cursfosc}{curs!72!black}%
  \def\fm@logo{\fm@logodir/vertical-#1.pdf}%
  \def\fm@marca{\fm@logodir/marca-#1.pdf}}

% ---------- capçalera de la primera pàgina ----------
% #1 matèria (petit), #2 curs, #3 tipus de document i unitat, #4 títol
\newcommand{\capcalera}[4]{%
  \gdef\fm@cursnom{#2}%
  \par\noindent
  \begin{tcolorbox}[colback=curs,colframe=curs,boxrule=0pt,arc=2.2mm,width=3.5cm,height=2.15cm,
      valign=center,left=3mm,right=2mm,top=1mm,bottom=1mm,nobeforeafter,box align=center]
    \color{curstext}{\footnotesize #1}\par\vspace{1pt}{\fontsize{19}{21}\selectfont\bfseries #2\par}%
  \end{tcolorbox}\hspace{5mm}%
  \begin{minipage}[c]{\dimexpr\linewidth-3.5cm-5mm-2.9cm\relax}
    {\color{cursfosc}\small\bfseries\MakeUppercase{#3}}\par\vspace{2pt}
    {\color{navy}\fontsize{20}{23}\selectfont\bfseries #4\par}
  \end{minipage}\hfill
  \begin{minipage}[c]{2.6cm}\raggedleft\includegraphics[width=2.6cm]{\fm@logo}\end{minipage}\par
  \vspace{2pt}}

% ---------- barres de secció del color del curs ----------
\renewcommand{\seccio}[2]{\par\needspace{\fmneed\baselineskip}\bigskip
  \begin{tcolorbox}[colback=curs,colframe=curs,coltext=curstext,arc=2mm,boxrule=0pt,left=3mm,right=3mm,top=1.2mm,bottom=1.2mm]
  {\large\bfseries #1\quad #2}
  \end{tcolorbox}\vspace{-2pt}}

% ---------- capçalera i peu de les pàgines ----------
% primera pàgina: logos institucionals al peu
\fancypagestyle{primera}{\fancyhf{}\renewcommand{\headrulewidth}{0pt}
  \fancyfoot[L]{\raisebox{-0.32\height}{\includegraphics[height=0.78cm]{logo-gva.png}}\hspace{5mm}%
    \raisebox{-0.32\height}{\includegraphics[height=0.78cm]{logo-vermellar.png}}}
  \fancyfoot[R]{\scriptsize\color{muted}\href{https://creativecommons.org/licenses/by-nc-sa/4.0/deed.ca}{\ccbyncsa}\ \ CC BY-NC-SA 4.0 · fentmates.org\qquad\small\thepage}}
% la resta: el símbol a l'esquerra i el curs a la dreta, amb el títol al centre com abans
\fancyhead[L]{\raisebox{-0.3\height}{\includegraphics[height=0.62cm]{\fm@marca}}}
\fancyhead[R]{\small\bfseries\color{cursfosc}\fm@cursnom}
\renewcommand{\headrule}{\color{curs!45}\hrule height \headrulewidth}
\makeatother

%   \fmcurs{1eso}
% i, al començament del document:
%   \capcalera{Matemàtiques}{1r ESO}{Fitxa d'exercicis · Unitat 2}{Divisibilitat}
% Opcional abans de \input: \def\fmneed{14} (espai mínim abans de cada secció).
% =====================================================================
\makeatletter
\def\fm@logodir{../../logo/pdf}
\providecommand{\fmneed}{10}
\def\fm@cursnom{}

% clau del curs, color i color del text damunt del color
\newcommand{\fm@curs}[3]{\@namedef{fm@color@#1}{#2}\@namedef{fm@text@#1}{#3}}
\fm@curs{1eso}{E6B84A}{navy}
\fm@curs{2eso}{3B82F6}{white}
\fm@curs{3eso}{FF7A50}{white}
\fm@curs{4eso-a}{A78BFA}{white}
\fm@curs{4eso-b}{6D28D9}{white}
\fm@curs{1batx-mi}{9B1C1C}{white}
\fm@curs{1batx-ccss}{EF6B6B}{white}
\fm@curs{2batx-mii}{1E6B47}{white}
\fm@curs{2batx-ccss}{2BB3A3}{white}

\newcommand{\fmcurs}[1]{%
  \@ifundefined{fm@color@#1}{\PackageError{fentmates}{Curs desconegut: #1}{}}{}%
  \edef\fm@tmp{\noexpand\definecolor{curs}{HTML}{\@nameuse{fm@color@#1}}%
               \noexpand\colorlet{curstext}{\@nameuse{fm@text@#1}}}\fm@tmp
  \colorlet{cursfosc}{curs!72!black}%
  \def\fm@logo{\fm@logodir/vertical-#1.pdf}%
  \def\fm@marca{\fm@logodir/marca-#1.pdf}}

% ---------- capçalera de la primera pàgina ----------
% #1 matèria (petit), #2 curs, #3 tipus de document i unitat, #4 títol
\newcommand{\capcalera}[4]{%
  \gdef\fm@cursnom{#2}%
  \par\noindent
  \begin{tcolorbox}[colback=curs,colframe=curs,boxrule=0pt,arc=2.2mm,width=3.5cm,height=2.15cm,
      valign=center,left=3mm,right=2mm,top=1mm,bottom=1mm,nobeforeafter,box align=center]
    \color{curstext}{\footnotesize #1}\par\vspace{1pt}{\fontsize{19}{21}\selectfont\bfseries #2\par}%
  \end{tcolorbox}\hspace{5mm}%
  \begin{minipage}[c]{\dimexpr\linewidth-3.5cm-5mm-2.9cm\relax}
    {\color{cursfosc}\small\bfseries\MakeUppercase{#3}}\par\vspace{2pt}
    {\color{navy}\fontsize{20}{23}\selectfont\bfseries #4\par}
  \end{minipage}\hfill
  \begin{minipage}[c]{2.6cm}\raggedleft\includegraphics[width=2.6cm]{\fm@logo}\end{minipage}\par
  \vspace{2pt}}

% ---------- barres de secció del color del curs ----------
\renewcommand{\seccio}[2]{\par\needspace{\fmneed\baselineskip}\bigskip
  \begin{tcolorbox}[colback=curs,colframe=curs,coltext=curstext,arc=2mm,boxrule=0pt,left=3mm,right=3mm,top=1.2mm,bottom=1.2mm]
  {\large\bfseries #1\quad #2}
  \end{tcolorbox}\vspace{-2pt}}

% ---------- capçalera i peu de les pàgines ----------
% primera pàgina: logos institucionals al peu
\fancypagestyle{primera}{\fancyhf{}\renewcommand{\headrulewidth}{0pt}
  \fancyfoot[L]{\raisebox{-0.32\height}{\includegraphics[height=0.78cm]{logo-gva.png}}\hspace{5mm}%
    \raisebox{-0.32\height}{\includegraphics[height=0.78cm]{logo-vermellar.png}}}
  \fancyfoot[R]{\scriptsize\color{muted}\href{https://creativecommons.org/licenses/by-nc-sa/4.0/deed.ca}{\ccbyncsa}\ \ CC BY-NC-SA 4.0 · fentmates.org\qquad\small\thepage}}
% la resta: el símbol a l'esquerra i el curs a la dreta, amb el títol al centre com abans
\fancyhead[L]{\raisebox{-0.3\height}{\includegraphics[height=0.62cm]{\fm@marca}}}
\fancyhead[R]{\small\bfseries\color{cursfosc}\fm@cursnom}
\renewcommand{\headrule}{\color{curs!45}\hrule height \headrulewidth}
\makeatother

%   \fmcurs{1eso}
% i, al començament del document:
%   \capcalera{Matemàtiques}{1r ESO}{Fitxa d'exercicis · Unitat 2}{Divisibilitat}
% Opcional abans de \input: \def\fmneed{14} (espai mínim abans de cada secció).
% =====================================================================
\makeatletter
\def\fm@logodir{../../logo/pdf}
\providecommand{\fmneed}{10}
\def\fm@cursnom{}

% clau del curs, color i color del text damunt del color
\newcommand{\fm@curs}[3]{\@namedef{fm@color@#1}{#2}\@namedef{fm@text@#1}{#3}}
\fm@curs{1eso}{E6B84A}{navy}
\fm@curs{2eso}{3B82F6}{white}
\fm@curs{3eso}{FF7A50}{white}
\fm@curs{4eso-a}{A78BFA}{white}
\fm@curs{4eso-b}{6D28D9}{white}
\fm@curs{1batx-mi}{9B1C1C}{white}
\fm@curs{1batx-ccss}{EF6B6B}{white}
\fm@curs{2batx-mii}{1E6B47}{white}
\fm@curs{2batx-ccss}{2BB3A3}{white}

\newcommand{\fmcurs}[1]{%
  \@ifundefined{fm@color@#1}{\PackageError{fentmates}{Curs desconegut: #1}{}}{}%
  \edef\fm@tmp{\noexpand\definecolor{curs}{HTML}{\@nameuse{fm@color@#1}}%
               \noexpand\colorlet{curstext}{\@nameuse{fm@text@#1}}}\fm@tmp
  \colorlet{cursfosc}{curs!72!black}%
  \def\fm@logo{\fm@logodir/vertical-#1.pdf}%
  \def\fm@marca{\fm@logodir/marca-#1.pdf}}

% ---------- capçalera de la primera pàgina ----------
% #1 matèria (petit), #2 curs, #3 tipus de document i unitat, #4 títol
\newcommand{\capcalera}[4]{%
  \gdef\fm@cursnom{#2}%
  \par\noindent
  \begin{tcolorbox}[colback=curs,colframe=curs,boxrule=0pt,arc=2.2mm,width=3.5cm,height=2.15cm,
      valign=center,left=3mm,right=2mm,top=1mm,bottom=1mm,nobeforeafter,box align=center]
    \color{curstext}{\footnotesize #1}\par\vspace{1pt}{\fontsize{19}{21}\selectfont\bfseries #2\par}%
  \end{tcolorbox}\hspace{5mm}%
  \begin{minipage}[c]{\dimexpr\linewidth-3.5cm-5mm-2.9cm\relax}
    {\color{cursfosc}\small\bfseries\MakeUppercase{#3}}\par\vspace{2pt}
    {\color{navy}\fontsize{20}{23}\selectfont\bfseries #4\par}
  \end{minipage}\hfill
  \begin{minipage}[c]{2.6cm}\raggedleft\includegraphics[width=2.6cm]{\fm@logo}\end{minipage}\par
  \vspace{2pt}}

% ---------- barres de secció del color del curs ----------
\renewcommand{\seccio}[2]{\par\needspace{\fmneed\baselineskip}\bigskip
  \begin{tcolorbox}[colback=curs,colframe=curs,coltext=curstext,arc=2mm,boxrule=0pt,left=3mm,right=3mm,top=1.2mm,bottom=1.2mm]
  {\large\bfseries #1\quad #2}
  \end{tcolorbox}\vspace{-2pt}}

% ---------- capçalera i peu de les pàgines ----------
% primera pàgina: logos institucionals al peu
\fancypagestyle{primera}{\fancyhf{}\renewcommand{\headrulewidth}{0pt}
  \fancyfoot[L]{\raisebox{-0.32\height}{\includegraphics[height=0.78cm]{logo-gva.png}}\hspace{5mm}%
    \raisebox{-0.32\height}{\includegraphics[height=0.78cm]{logo-vermellar.png}}}
  \fancyfoot[R]{\scriptsize\color{muted}\href{https://creativecommons.org/licenses/by-nc-sa/4.0/deed.ca}{\ccbyncsa}\ \ CC BY-NC-SA 4.0 · fentmates.org\qquad\small\thepage}}
% la resta: el símbol a l'esquerra i el curs a la dreta, amb el títol al centre com abans
\fancyhead[L]{\raisebox{-0.3\height}{\includegraphics[height=0.62cm]{\fm@marca}}}
\fancyhead[R]{\small\bfseries\color{cursfosc}\fm@cursnom}
\renewcommand{\headrule}{\color{curs!45}\hrule height \headrulewidth}
\makeatother

\fmcurs{3eso}

\begin{document}
\thispagestyle{primera}

\capcalera{Matemàtiques}{3r ESO}{Fitxa d'exercicis · Unitat 7}{Funcions}
\vspace{10pt}
Nom i cognoms: \rule{8.4cm}{0.4pt}\hfill Grup: \rule{1.6cm}{0.4pt}\hfill Data: \rule{2.2cm}{0.4pt}
\vspace{8pt}

\begin{tcolorbox}[colback=soft,colframe=soft,arc=2mm,left=3mm,right=3mm,top=1.5mm,bottom=1.5mm]
\small Exercicis graduats per nivells de dificultat:\quad \estrelles{1}\ bàsic \qquad \estrelles{2}\ mitjà \qquad \estrelles{3}\ avançat.\\[2pt] Resol-los al quadern mostrant tots els passos. Per a dibuixar gràfiques, fes servir paper quadriculat i tria una escala adequada per a cada eix.
\end{tcolorbox}

% =====================================================
\seccio{1}{Concepte de funció}

\ex{1} Indica si cada relació és una funció:\begin{tasks}(2)
  \task A cada nombre natural, el seu quadrat.
  \task A cada persona, els seus germans.
  \task A cada quantitat de litres de gasolina, el preu que paguem.
  \task A cada $x>0$, les seues arrels quadrades.
  \task A cada alumne, la nota de l'últim examen.
  \task A cada nota, els alumnes que l'han tret.
\end{tasks}\recorda{Una relació és una funció si a \textbf{cada} valor de la variable independent ($x$) li correspon \textbf{un únic} valor de la variable dependent ($y$).}

\ex{1} En cada relació, quina és la variable independent i quina la dependent?\begin{tasks}(1)
  \task La durada d'una trucada i el que costa.
  \task El nombre d'obrers i el temps que tarden a acabar una obra.
  \task El costat d'un triangle equilàter i el seu perímetre.
\end{tasks}

\ex{1} Quines d'aquestes gràfiques corresponen a una funció?\par\smallskip\begin{tasks}(4)
  \task \begin{tikzpicture}[x=0.27cm,y=0.27cm,baseline=(current bounding box.center)]
\draw[muted!22,very thin,xstep=1,ystep=1] (-5,-4) grid (5,4);
\draw[-latex,muted] (-5,0)--(5.35,0) node[right,font=\scriptsize]{$x$};
\draw[-latex,muted] (0,-4)--(0,4.35) node[above,font=\scriptsize]{$y$};
\node[font=\tiny,muted,below] at (-4,0) {-4};
\node[font=\tiny,muted,below] at (-2,0) {-2};
\node[font=\tiny,muted,below] at (2,0) {2};
\node[font=\tiny,muted,below] at (4,0) {4};
\node[font=\tiny,muted,left] at (0,-4) {-4};
\node[font=\tiny,muted,left] at (0,-2) {-2};
\node[font=\tiny,muted,left] at (0,2) {2};
\begin{scope}\clip (-5,-4) rectangle (5,4);
\draw[accent,line width=1.1pt] (-4.0,2.0)--(-2.0,-2.0)--(1.0,3.0)--(4.0,1.0);
\end{scope}
\fill[accent] (-4.0,2.0) circle (2.2pt);
\fill[accent] (4.0,1.0) circle (2.2pt);
\end{tikzpicture}
  \task \begin{tikzpicture}[x=0.27cm,y=0.27cm,baseline=(current bounding box.center)]
\draw[muted!22,very thin,xstep=1,ystep=1] (-5,-4) grid (5,4);
\draw[-latex,muted] (-5,0)--(5.35,0) node[right,font=\scriptsize]{$x$};
\draw[-latex,muted] (0,-4)--(0,4.35) node[above,font=\scriptsize]{$y$};
\node[font=\tiny,muted,below] at (-4,0) {-4};
\node[font=\tiny,muted,below] at (-2,0) {-2};
\node[font=\tiny,muted,below] at (2,0) {2};
\node[font=\tiny,muted,below] at (4,0) {4};
\node[font=\tiny,muted,left] at (0,-4) {-4};
\node[font=\tiny,muted,left] at (0,-2) {-2};
\node[font=\tiny,muted,left] at (0,2) {2};
\begin{scope}\clip (-5,-4) rectangle (5,4);
\draw[teal,line width=1.1pt] plot[smooth] coordinates {(3.000,0.000) (2.984,0.314) (2.934,0.624) (2.853,0.927) (2.741,1.220) (2.598,1.500) (2.427,1.763) (2.229,2.007) (2.007,2.229) (1.763,2.427) (1.500,2.598) (1.220,2.741) (0.927,2.853) (0.624,2.934) (0.314,2.984) (0.000,3.000) (-0.314,2.984) (-0.624,2.934) (-0.927,2.853) (-1.220,2.741) (-1.500,2.598) (-1.763,2.427) (-2.007,2.229) (-2.229,2.007) (-2.427,1.763) (-2.598,1.500) (-2.741,1.220) (-2.853,0.927) (-2.934,0.624) (-2.984,0.314) (-3.000,0.000) (-2.984,-0.314) (-2.934,-0.624) (-2.853,-0.927) (-2.741,-1.220) (-2.598,-1.500) (-2.427,-1.763) (-2.229,-2.007) (-2.007,-2.229) (-1.763,-2.427) (-1.500,-2.598) (-1.220,-2.741) (-0.927,-2.853) (-0.624,-2.934) (-0.314,-2.984) (-0.000,-3.000) (0.314,-2.984) (0.624,-2.934) (0.927,-2.853) (1.220,-2.741) (1.500,-2.598) (1.763,-2.427) (2.007,-2.229) (2.229,-2.007) (2.427,-1.763) (2.598,-1.500) (2.741,-1.220) (2.853,-0.927) (2.934,-0.624) (2.984,-0.314) (3.000,-0.000)};
\end{scope}
\end{tikzpicture}
  \task \begin{tikzpicture}[x=0.27cm,y=0.27cm,baseline=(current bounding box.center)]
\draw[muted!22,very thin,xstep=1,ystep=1] (-5,-4) grid (5,4);
\draw[-latex,muted] (-5,0)--(5.35,0) node[right,font=\scriptsize]{$x$};
\draw[-latex,muted] (0,-4)--(0,4.35) node[above,font=\scriptsize]{$y$};
\node[font=\tiny,muted,below] at (-4,0) {-4};
\node[font=\tiny,muted,below] at (-2,0) {-2};
\node[font=\tiny,muted,below] at (2,0) {2};
\node[font=\tiny,muted,below] at (4,0) {4};
\node[font=\tiny,muted,left] at (0,-4) {-4};
\node[font=\tiny,muted,left] at (0,-2) {-2};
\node[font=\tiny,muted,left] at (0,2) {2};
\begin{scope}\clip (-5,-4) rectangle (5,4);
\draw[teal,line width=1.1pt] plot[smooth] coordinates {(5.000,-4.000) (4.605,-3.900) (4.220,-3.800) (3.845,-3.700) (3.480,-3.600) (3.125,-3.500) (2.780,-3.400) (2.445,-3.300) (2.120,-3.200) (1.805,-3.100) (1.500,-3.000) (1.205,-2.900) (0.920,-2.800) (0.645,-2.700) (0.380,-2.600) (0.125,-2.500) (-0.120,-2.400) (-0.355,-2.300) (-0.580,-2.200) (-0.795,-2.100) (-1.000,-2.000) (-1.195,-1.900) (-1.380,-1.800) (-1.555,-1.700) (-1.720,-1.600) (-1.875,-1.500) (-2.020,-1.400) (-2.155,-1.300) (-2.280,-1.200) (-2.395,-1.100) (-2.500,-1.000) (-2.595,-0.900) (-2.680,-0.800) (-2.755,-0.700) (-2.820,-0.600) (-2.875,-0.500) (-2.920,-0.400) (-2.955,-0.300) (-2.980,-0.200) (-2.995,-0.100) (-3.000,0.000) (-2.995,0.100) (-2.980,0.200) (-2.955,0.300) (-2.920,0.400) (-2.875,0.500) (-2.820,0.600) (-2.755,0.700) (-2.680,0.800) (-2.595,0.900) (-2.500,1.000) (-2.395,1.100) (-2.280,1.200) (-2.155,1.300) (-2.020,1.400) (-1.875,1.500) (-1.720,1.600) (-1.555,1.700) (-1.380,1.800) (-1.195,1.900) (-1.000,2.000) (-0.795,2.100) (-0.580,2.200) (-0.355,2.300) (-0.120,2.400) (0.125,2.500) (0.380,2.600) (0.645,2.700) (0.920,2.800) (1.205,2.900) (1.500,3.000) (1.805,3.100) (2.120,3.200) (2.445,3.300) (2.780,3.400) (3.125,3.500) (3.480,3.600) (3.845,3.700) (4.220,3.800) (4.605,3.900) (5.000,4.000)};
\end{scope}
\end{tikzpicture}
  \task \begin{tikzpicture}[x=0.27cm,y=0.27cm,baseline=(current bounding box.center)]
\draw[muted!22,very thin,xstep=1,ystep=1] (-5,-4) grid (5,4);
\draw[-latex,muted] (-5,0)--(5.35,0) node[right,font=\scriptsize]{$x$};
\draw[-latex,muted] (0,-4)--(0,4.35) node[above,font=\scriptsize]{$y$};
\node[font=\tiny,muted,below] at (-4,0) {-4};
\node[font=\tiny,muted,below] at (-2,0) {-2};
\node[font=\tiny,muted,below] at (2,0) {2};
\node[font=\tiny,muted,below] at (4,0) {4};
\node[font=\tiny,muted,left] at (0,-4) {-4};
\node[font=\tiny,muted,left] at (0,-2) {-2};
\node[font=\tiny,muted,left] at (0,2) {2};
\begin{scope}\clip (-5,-4) rectangle (5,4);
\draw[teal,line width=1.1pt] (-4.000,-3.000)--(0.000,3.000)--(-2.000,0.000)--(4.000,2.000);
\end{scope}
\end{tikzpicture}
\end{tasks}

\ex{2} Les taules següents, defineixen una funció de $x$?\begin{tasks}(2)
  \task $\begin{array}{|c|c|c|c|c|}\hline x&1&2&2&3\\\hline y&4&5&6&7\\\hline\end{array}$
  \task $\begin{array}{|c|c|c|c|c|}\hline x&1&2&3&4\\\hline y&5&5&5&6\\\hline\end{array}$
\end{tasks}

\ex{2} Explica per què la recta $y=3$ és la gràfica d'una funció i la recta $x=3$ no ho és.

\ex{3} Dibuixa una gràfica que passe pels punts $(-2,\,1)$, $(0,\,3)$ i $(3,\,1)$ i que siga una funció, i una altra que passe pels mateixos punts i que \textbf{no} ho siga.

\ex{3} A cada nombre $x$ li fem correspondre el nombre $y$ que compleix $y^2=x$. És una funció? I si la condició és $y^3=x$?


% =====================================================
\seccio{2}{Expressió d'una funció}

\ex{1} Donada la funció $f(x)=3x-5$, calcula:\begin{tasks}(4)
  \task $f(0)$
  \task $f(2)$
  \task $f(-3)$
  \task $f\!\left(\tfrac12\right)$
\end{tasks}\recorda{Per a calcular la imatge $f(a)$, substituïm la $x$ per $a$ en la fórmula. Els nombres negatius, entre parèntesis.}

\ex{1} Si $g(x)=x^2-4x+3$, calcula $g(0)$, $g(1)$, $g(-2)$ i $g(5)$.

\ex{1} Completa les taules de valors:\begin{tasks}(2)
  \task $y=2x-1$\\[3pt]$\begin{array}{|c|C{0.7cm}C{0.7cm}C{0.7cm}C{0.7cm}C{0.7cm}|}\hline x&-2&-1&0&1&2\\\hline y&\phantom{00}&\phantom{00}&\phantom{00}&\phantom{00}&\phantom{00}\\\hline\end{array}$
  \task $y=-x^2+3$\\[3pt]$\begin{array}{|c|C{0.7cm}C{0.7cm}C{0.7cm}C{0.7cm}C{0.7cm}|}\hline x&-2&-1&0&1&2\\\hline y&\phantom{00}&\phantom{00}&\phantom{00}&\phantom{00}&\phantom{00}\\\hline\end{array}$
\end{tasks}

\ex{2} Calcula les imatges que s'indiquen:\begin{tasks}(2)
  \task $h(x)=\dfrac{2x+1}{3}$: \ $h(4)$, $h(-2)$, $h(1)$
  \task $k(x)=\sqrt{x+5}$: \ $k(4)$, $k(-1)$, $k(11)$
\end{tasks}

\ex{2} Donada $f(x)=5x-3$, troba l'antiimatge de:\begin{tasks}(3)
  \task $12$
  \task $-8$
  \task $0$
\end{tasks}\recorda{Per a trobar l'antiimatge de $k$, resolem l'equació $f(x)=k$.}

\ex{2} Donada $f(x)=x^2-1$, quins valors de $x$ tenen per imatge:\begin{tasks}(3)
  \task $8$
  \task $-1$
  \task $-5$
\end{tasks}

\ex{2} Escriu la fórmula de cada funció:\begin{tasks}(1)
  \task El perímetre d'un rectangle de base $x$ i altura 4 cm.
  \task El preu de $x$ entrades de concert a 8,50 € cadascuna.
  \task L'àrea d'un cercle de radi $x$.
  \task El temps que es tarda a fer 300 km a una velocitat de $x$ km/h.
\end{tasks}

\ex{2} Troba la fórmula de les funcions que donen aquestes taules:\begin{tasks}(2)
  \task $\begin{array}{|c|ccccc|}\hline x&0&1&2&3&4\\\hline y&3&7&11&15&19\\\hline\end{array}$
  \task $\begin{array}{|c|cccc|}\hline x&1&2&3&4\\\hline y&1&4&9&16\\\hline\end{array}$
\end{tasks}

\ex{3} Una botiga de fotografia cobra 2 € per l'enviament més 0,15 € per cada foto revelada. Escriu la fórmula del preu en funció del nombre de fotos, calcula el preu de 40 fotos i quantes fotos podem revelar amb 11 €.

\ex{3} La funció $f(x)=ax+3$ compleix que $f(2)=11$. Troba $a$ i calcula $f(-1)$.


% =====================================================
\seccio{3}{Les funcions i els seus gràfics}

\ex{1} Fes una taula de valors de $f(x)=2x-3$ per a $x=-1,\,0,\,1,\,2,\,3$, representa la funció i troba els punts de tall amb els eixos.\recorda{Tall amb l'eix $Y$: fem $x=0$. Tall amb l'eix $X$: resolem $f(x)=0$.}

\ex{1} Fes una taula de valors de $g(x)=x^2-4$ per a $x$ entre $-3$ i $3$ i representa-la. On talla els eixos?

\ex{2} Troba els punts de tall amb els eixos:\begin{tasks}(2)
  \task $f(x)=3x+6$
  \task $g(x)=x^2-5x+6$
  \task $h(x)=x^2+1$
  \task $k(x)=x^3-4x$
\end{tasks}

\ex{2} La gràfica mostra la distància a casa d'un ciclista durant una eixida.\par\smallskip\begin{center}\begin{tikzpicture}[x=0.9cm,y=0.07cm,baseline=(current bounding box.center)]
\draw[muted!22,very thin,xstep=1,ystep=5] (0,0) grid (7,45);
\draw[-latex,muted] (0,0)--(7.35,0) node[right,font=\scriptsize]{$t\ (\mathrm{h})$};
\draw[-latex,muted] (0,0)--(0,46.75) node[above,font=\scriptsize]{$d\ (\mathrm{km})$};
\node[font=\tiny,muted,below] at (1,0) {1};
\node[font=\tiny,muted,below] at (2,0) {2};
\node[font=\tiny,muted,below] at (3,0) {3};
\node[font=\tiny,muted,below] at (4,0) {4};
\node[font=\tiny,muted,below] at (5,0) {5};
\node[font=\tiny,muted,below] at (6,0) {6};
\node[font=\tiny,muted,left] at (0,10) {10};
\node[font=\tiny,muted,left] at (0,20) {20};
\node[font=\tiny,muted,left] at (0,30) {30};
\node[font=\tiny,muted,left] at (0,40) {40};
\begin{scope}\clip (0,0) rectangle (7,45);
\draw[accent,line width=1.1pt] (0.0,0.0)--(2.0,30.0)--(3.0,30.0)--(4.0,40.0)--(6.0,0.0);
\end{scope}
\fill[accent] (0.0,0.0) circle (2.2pt);
\fill[accent] (6.0,0.0) circle (2.2pt);
\end{tikzpicture}\end{center}\begin{tasks}(1)
  \task Quants quilòmetres recorre en total?
  \task Quant de temps està parat?
  \task A quina velocitat va en les dues primeres hores? I en la tornada?
  \task A quina distància de casa està a les 5 h?
  \task En quins moments està a 15 km de casa?
\end{tasks}

\ex{2}\label{ex:temp} La gràfica mostra la temperatura d'un dia d'hivern en una ciutat.\par\smallskip\begin{center}\begin{tikzpicture}[x=0.3cm,y=0.17cm,baseline=(current bounding box.center)]
\draw[muted!22,very thin,xstep=1,ystep=1] (0,-2) grid (25,18);
\draw[-latex,muted] (0,0)--(25.35,0) node[right,font=\scriptsize]{$h$};
\draw[-latex,muted] (0,-2)--(0,18.35) node[above,font=\scriptsize]{$T\ (^\circ\mathrm{C})$};
\node[font=\tiny,muted,below] at (2,0) {2};
\node[font=\tiny,muted,below] at (4,0) {4};
\node[font=\tiny,muted,below] at (6,0) {6};
\node[font=\tiny,muted,below] at (8,0) {8};
\node[font=\tiny,muted,below] at (10,0) {10};
\node[font=\tiny,muted,below] at (12,0) {12};
\node[font=\tiny,muted,below] at (14,0) {14};
\node[font=\tiny,muted,below] at (16,0) {16};
\node[font=\tiny,muted,below] at (18,0) {18};
\node[font=\tiny,muted,below] at (20,0) {20};
\node[font=\tiny,muted,below] at (22,0) {22};
\node[font=\tiny,muted,below] at (24,0) {24};
\node[font=\tiny,muted,left] at (0,-2) {-2};
\node[font=\tiny,muted,left] at (0,2) {2};
\node[font=\tiny,muted,left] at (0,4) {4};
\node[font=\tiny,muted,left] at (0,6) {6};
\node[font=\tiny,muted,left] at (0,8) {8};
\node[font=\tiny,muted,left] at (0,10) {10};
\node[font=\tiny,muted,left] at (0,12) {12};
\node[font=\tiny,muted,left] at (0,14) {14};
\node[font=\tiny,muted,left] at (0,16) {16};
\begin{scope}\clip (0,-2) rectangle (25,18);
\draw[accent,line width=1.1pt] (0.0,4.0)--(6.0,0.0)--(14.0,16.0)--(20.0,10.0)--(24.0,6.0);
\end{scope}
\fill[accent] (0.0,4.0) circle (2.2pt);
\fill[accent] (24.0,6.0) circle (2.2pt);
\end{tikzpicture}\end{center}\begin{tasks}(2)
  \task Quina temperatura hi havia a les 10 h?
  \task A quina hora va fer més calor? I més fred?
  \task Quan hi havia 10 °C?
  \task Quant va pujar la temperatura de 6 h a 14 h?
\end{tasks}

\ex{2} Associa cada situació amb la seua gràfica:\begin{tasks}(1)
  \task El preu de la fruita segons els quilos que en comprem.
  \task L'altura d'una pilota que llancem cap amunt.
  \task La temperatura d'un cafè calent que es refreda.
  \task L'aigua d'una banyera que omplim, deixem plena una estona i després buidem.
\end{tasks}\par\smallskip\begin{center}\textbf{I}\ \begin{tikzpicture}[x=0.3cm,y=0.3cm,baseline=(current bounding box.center)]
\draw[muted!22,very thin,xstep=1,ystep=1] (-1,-1) grid (6,6);
\draw[-latex,muted] (-1,0)--(6.35,0) node[right,font=\scriptsize]{$t$};
\draw[-latex,muted] (0,-1)--(0,6.35) node[above,font=\scriptsize]{$y$};
\begin{scope}\clip (-1,-1) rectangle (6,6);
\draw[accent,line width=1.1pt] plot[smooth] coordinates {(0.000,0.000) (0.031,0.124) (0.062,0.247) (0.094,0.368) (0.125,0.487) (0.156,0.605) (0.188,0.722) (0.219,0.837) (0.250,0.950) (0.281,1.062) (0.312,1.172) (0.344,1.280) (0.375,1.387) (0.406,1.493) (0.438,1.597) (0.469,1.699) (0.500,1.800) (0.531,1.899) (0.562,1.997) (0.594,2.093) (0.625,2.188) (0.656,2.280) (0.688,2.372) (0.719,2.462) (0.750,2.550) (0.781,2.637) (0.812,2.722) (0.844,2.805) (0.875,2.887) (0.906,2.968) (0.938,3.047) (0.969,3.124) (1.000,3.200) (1.031,3.274) (1.062,3.347) (1.094,3.418) (1.125,3.487) (1.156,3.555) (1.188,3.622) (1.219,3.687) (1.250,3.750) (1.281,3.812) (1.312,3.872) (1.344,3.930) (1.375,3.987) (1.406,4.043) (1.438,4.097) (1.469,4.149) (1.500,4.200) (1.531,4.249) (1.562,4.297) (1.594,4.343) (1.625,4.388) (1.656,4.430) (1.688,4.472) (1.719,4.512) (1.750,4.550) (1.781,4.587) (1.812,4.622) (1.844,4.655) (1.875,4.688) (1.906,4.718) (1.938,4.747) (1.969,4.774) (2.000,4.800) (2.031,4.824) (2.062,4.847) (2.094,4.868) (2.125,4.888) (2.156,4.905) (2.188,4.922) (2.219,4.937) (2.250,4.950) (2.281,4.962) (2.312,4.972) (2.344,4.980) (2.375,4.987) (2.406,4.993) (2.438,4.997) (2.469,4.999) (2.500,5.000) (2.531,4.999) (2.562,4.997) (2.594,4.993) (2.625,4.987) (2.656,4.980) (2.688,4.972) (2.719,4.962) (2.750,4.950) (2.781,4.937) (2.812,4.922) (2.844,4.905) (2.875,4.888) (2.906,4.868) (2.938,4.847) (2.969,4.824) (3.000,4.800) (3.031,4.774) (3.062,4.747) (3.094,4.718) (3.125,4.688) (3.156,4.655) (3.188,4.622) (3.219,4.587) (3.250,4.550) (3.281,4.512) (3.312,4.472) (3.344,4.430) (3.375,4.388) (3.406,4.343) (3.438,4.297) (3.469,4.249) (3.500,4.200) (3.531,4.149) (3.562,4.097) (3.594,4.043) (3.625,3.987) (3.656,3.930) (3.688,3.872) (3.719,3.812) (3.750,3.750) (3.781,3.687) (3.812,3.622) (3.844,3.555) (3.875,3.487) (3.906,3.418) (3.938,3.347) (3.969,3.274) (4.000,3.200) (4.031,3.124) (4.062,3.047) (4.094,2.968) (4.125,2.887) (4.156,2.805) (4.188,2.722) (4.219,2.637) (4.250,2.550) (4.281,2.462) (4.312,2.372) (4.344,2.280) (4.375,2.188) (4.406,2.093) (4.438,1.997) (4.469,1.899) (4.500,1.800) (4.531,1.699) (4.562,1.597) (4.594,1.493) (4.625,1.387) (4.656,1.280) (4.688,1.172) (4.719,1.062) (4.750,0.950) (4.781,0.837) (4.812,0.722) (4.844,0.605) (4.875,0.487) (4.906,0.368) (4.938,0.247) (4.969,0.124) (5.000,0.000)};
\end{scope}
\end{tikzpicture}\hfill\textbf{II}\ \begin{tikzpicture}[x=0.3cm,y=0.3cm,baseline=(current bounding box.center)]
\draw[muted!22,very thin,xstep=1,ystep=1] (-1,-1) grid (6,6);
\draw[-latex,muted] (-1,0)--(6.35,0) node[right,font=\scriptsize]{$t$};
\draw[-latex,muted] (0,-1)--(0,6.35) node[above,font=\scriptsize]{$y$};
\begin{scope}\clip (-1,-1) rectangle (6,6);
\draw[accent,line width=1.1pt] (0.0,0.0)--(2.0,5.0)--(4.0,5.0)--(5.0,0.0);
\end{scope}
\fill[accent] (0.0,0.0) circle (2.2pt);
\fill[accent] (5.0,0.0) circle (2.2pt);
\end{tikzpicture}\hfill\textbf{III}\ \begin{tikzpicture}[x=0.3cm,y=0.3cm,baseline=(current bounding box.center)]
\draw[muted!22,very thin,xstep=1,ystep=1] (-1,-1) grid (6,6);
\draw[-latex,muted] (-1,0)--(6.35,0) node[right,font=\scriptsize]{$x$};
\draw[-latex,muted] (0,-1)--(0,6.35) node[above,font=\scriptsize]{$y$};
\begin{scope}\clip (-1,-1) rectangle (6,6);
\draw[accent,line width=1.1pt] plot[smooth] coordinates {(0.000,0.000) (0.034,0.034) (0.069,0.069) (0.103,0.103) (0.138,0.138) (0.172,0.172) (0.206,0.206) (0.241,0.241) (0.275,0.275) (0.309,0.309) (0.344,0.344) (0.378,0.378) (0.412,0.412) (0.447,0.447) (0.481,0.481) (0.516,0.516) (0.550,0.550) (0.584,0.584) (0.619,0.619) (0.653,0.653) (0.688,0.688) (0.722,0.722) (0.756,0.756) (0.791,0.791) (0.825,0.825) (0.859,0.859) (0.894,0.894) (0.928,0.928) (0.963,0.963) (0.997,0.997) (1.031,1.031) (1.066,1.066) (1.100,1.100) (1.134,1.134) (1.169,1.169) (1.203,1.203) (1.238,1.238) (1.272,1.272) (1.306,1.306) (1.341,1.341) (1.375,1.375) (1.409,1.409) (1.444,1.444) (1.478,1.478) (1.512,1.512) (1.547,1.547) (1.581,1.581) (1.616,1.616) (1.650,1.650) (1.684,1.684) (1.719,1.719) (1.753,1.753) (1.788,1.788) (1.822,1.822) (1.856,1.856) (1.891,1.891) (1.925,1.925) (1.959,1.959) (1.994,1.994) (2.028,2.028) (2.062,2.062) (2.097,2.097) (2.131,2.131) (2.166,2.166) (2.200,2.200) (2.234,2.234) (2.269,2.269) (2.303,2.303) (2.337,2.337) (2.372,2.372) (2.406,2.406) (2.441,2.441) (2.475,2.475) (2.509,2.509) (2.544,2.544) (2.578,2.578) (2.612,2.612) (2.647,2.647) (2.681,2.681) (2.716,2.716) (2.750,2.750) (2.784,2.784) (2.819,2.819) (2.853,2.853) (2.888,2.888) (2.922,2.922) (2.956,2.956) (2.991,2.991) (3.025,3.025) (3.059,3.059) (3.094,3.094) (3.128,3.128) (3.163,3.163) (3.197,3.197) (3.231,3.231) (3.266,3.266) (3.300,3.300) (3.334,3.334) (3.369,3.369) (3.403,3.403) (3.438,3.438) (3.472,3.472) (3.506,3.506) (3.541,3.541) (3.575,3.575) (3.609,3.609) (3.644,3.644) (3.678,3.678) (3.712,3.712) (3.747,3.747) (3.781,3.781) (3.816,3.816) (3.850,3.850) (3.884,3.884) (3.919,3.919) (3.953,3.953) (3.987,3.987) (4.022,4.022) (4.056,4.056) (4.091,4.091) (4.125,4.125) (4.159,4.159) (4.194,4.194) (4.228,4.228) (4.263,4.263) (4.297,4.297) (4.331,4.331) (4.366,4.366) (4.400,4.400) (4.434,4.434) (4.469,4.469) (4.503,4.503) (4.537,4.537) (4.572,4.572) (4.606,4.606) (4.641,4.641) (4.675,4.675) (4.709,4.709) (4.744,4.744) (4.778,4.778) (4.812,4.812) (4.847,4.847) (4.881,4.881) (4.916,4.916) (4.950,4.950) (4.984,4.984) (5.019,5.019) (5.053,5.053) (5.088,5.088) (5.122,5.122) (5.156,5.156) (5.191,5.191) (5.225,5.225) (5.259,5.259) (5.294,5.294) (5.328,5.328) (5.362,5.362) (5.397,5.397) (5.431,5.431) (5.466,5.466) (5.500,5.500)};
\end{scope}
\end{tikzpicture}\hfill\textbf{IV}\ \begin{tikzpicture}[x=0.3cm,y=0.3cm,baseline=(current bounding box.center)]
\draw[muted!22,very thin,xstep=1,ystep=1] (-1,-1) grid (6,6);
\draw[-latex,muted] (-1,0)--(6.35,0) node[right,font=\scriptsize]{$t$};
\draw[-latex,muted] (0,-1)--(0,6.35) node[above,font=\scriptsize]{$y$};
\begin{scope}\clip (-1,-1) rectangle (6,6);
\draw[accent,line width=1.1pt] plot[smooth] coordinates {(0.000,5.000) (0.034,4.918) (0.069,4.838) (0.103,4.760) (0.138,4.683) (0.172,4.608) (0.206,4.534) (0.241,4.462) (0.275,4.392) (0.309,4.322) (0.344,4.255) (0.378,4.188) (0.412,4.123) (0.447,4.059) (0.481,3.997) (0.516,3.936) (0.550,3.876) (0.584,3.817) (0.619,3.759) (0.653,3.703) (0.688,3.648) (0.722,3.594) (0.756,3.541) (0.791,3.489) (0.825,3.438) (0.859,3.389) (0.894,3.340) (0.928,3.292) (0.963,3.245) (0.997,3.199) (1.031,3.154) (1.066,3.110) (1.100,3.067) (1.134,3.025) (1.169,2.984) (1.203,2.943) (1.238,2.904) (1.272,2.865) (1.306,2.827) (1.341,2.789) (1.375,2.753) (1.409,2.717) (1.444,2.682) (1.478,2.648) (1.512,2.614) (1.547,2.581) (1.581,2.549) (1.616,2.517) (1.650,2.486) (1.684,2.456) (1.719,2.426) (1.753,2.397) (1.788,2.369) (1.822,2.341) (1.856,2.313) (1.891,2.286) (1.925,2.260) (1.959,2.235) (1.994,2.209) (2.028,2.185) (2.062,2.160) (2.097,2.137) (2.131,2.114) (2.166,2.091) (2.200,2.069) (2.234,2.047) (2.269,2.025) (2.303,2.004) (2.337,1.984) (2.372,1.964) (2.406,1.944) (2.441,1.925) (2.475,1.906) (2.509,1.888) (2.544,1.869) (2.578,1.852) (2.612,1.834) (2.647,1.817) (2.681,1.801) (2.716,1.784) (2.750,1.768) (2.784,1.753) (2.819,1.737) (2.853,1.722) (2.888,1.707) (2.922,1.693) (2.956,1.679) (2.991,1.665) (3.025,1.651) (3.059,1.638) (3.094,1.625) (3.128,1.612) (3.163,1.600) (3.197,1.588) (3.231,1.576) (3.266,1.564) (3.300,1.552) (3.334,1.541) (3.369,1.530) (3.403,1.519) (3.438,1.509) (3.472,1.498) (3.506,1.488) (3.541,1.478) (3.575,1.468) (3.609,1.459) (3.644,1.449) (3.678,1.440) (3.712,1.431) (3.747,1.422) (3.781,1.414) (3.816,1.405) (3.850,1.397) (3.884,1.389) (3.919,1.381) (3.953,1.373) (3.987,1.366) (4.022,1.358) (4.056,1.351) (4.091,1.344) (4.125,1.337) (4.159,1.330) (4.194,1.323) (4.228,1.316) (4.263,1.310) (4.297,1.304) (4.331,1.297) (4.366,1.291) (4.400,1.285) (4.434,1.280) (4.469,1.274) (4.503,1.268) (4.537,1.263) (4.572,1.257) (4.606,1.252) (4.641,1.247) (4.675,1.242) (4.709,1.237) (4.744,1.232) (4.778,1.228) (4.812,1.223) (4.847,1.218) (4.881,1.214) (4.916,1.209) (4.950,1.205) (4.984,1.201) (5.019,1.197) (5.053,1.193) (5.088,1.189) (5.122,1.185) (5.156,1.181) (5.191,1.178) (5.225,1.174) (5.259,1.170) (5.294,1.167) (5.328,1.164) (5.362,1.160) (5.397,1.157) (5.431,1.154) (5.466,1.151) (5.500,1.148)};
\end{scope}
\end{tikzpicture}\end{center}

\ex{3} Fes una taula de valors de $f(x)=\dfrac{12}{x}$ per a $x=\pm1,\,\pm2,\,\pm3,\,\pm4,\,\pm6$ i representa-la. Per què no podem unir el punt de $x=-1$ amb el de $x=1$?

\ex{3} Inventa una història que corresponga a aquesta gràfica de la distància a casa d'una persona.\par\smallskip\begin{center}\begin{tikzpicture}[x=0.8cm,y=0.5cm,baseline=(current bounding box.center)]
\draw[muted!22,very thin,xstep=1,ystep=1] (0,0) grid (6,4);
\draw[-latex,muted] (0,0)--(6.35,0) node[right,font=\scriptsize]{$t\ (\mathrm{h})$};
\draw[-latex,muted] (0,0)--(0,4.35) node[above,font=\scriptsize]{$d\ (\mathrm{km})$};
\node[font=\tiny,muted,below] at (1,0) {1};
\node[font=\tiny,muted,below] at (2,0) {2};
\node[font=\tiny,muted,below] at (3,0) {3};
\node[font=\tiny,muted,below] at (4,0) {4};
\node[font=\tiny,muted,below] at (5,0) {5};
\node[font=\tiny,muted,left] at (0,1) {1};
\node[font=\tiny,muted,left] at (0,2) {2};
\node[font=\tiny,muted,left] at (0,3) {3};
\begin{scope}\clip (0,0) rectangle (6,4);
\draw[accent,line width=1.1pt] (0.0,0.0)--(1.0,2.0)--(3.0,2.0)--(3.5,0.0)--(4.0,0.0)--(5.0,3.0);
\end{scope}
\fill[accent] (0.0,0.0) circle (2.2pt);
\fill[accent] (5.0,3.0) circle (2.2pt);
\end{tikzpicture}\end{center}


% =====================================================
\seccio{4}{Domini i recorregut}

\ex{1} Escriu el domini i el recorregut d'aquestes funcions:\begin{tasks}(3)
  \task \begin{tikzpicture}[x=0.27cm,y=0.27cm,baseline=(current bounding box.center)]
\draw[muted!22,very thin,xstep=1,ystep=1] (-5,-3) grid (5,4);
\draw[-latex,muted] (-5,0)--(5.35,0) node[right,font=\scriptsize]{$x$};
\draw[-latex,muted] (0,-3)--(0,4.35) node[above,font=\scriptsize]{$y$};
\node[font=\tiny,muted,below] at (-4,0) {-4};
\node[font=\tiny,muted,below] at (-2,0) {-2};
\node[font=\tiny,muted,below] at (2,0) {2};
\node[font=\tiny,muted,below] at (4,0) {4};
\node[font=\tiny,muted,left] at (0,-2) {-2};
\node[font=\tiny,muted,left] at (0,2) {2};
\begin{scope}\clip (-5,-3) rectangle (5,4);
\draw[accent,line width=1.1pt] (-4.0,3.0)--(-1.0,-2.0)--(2.0,1.0)--(4.0,1.0);
\end{scope}
\fill[accent] (-4.0,3.0) circle (2.2pt);
\fill[accent] (4.0,1.0) circle (2.2pt);
\end{tikzpicture}
  \task \begin{tikzpicture}[x=0.27cm,y=0.27cm,baseline=(current bounding box.center)]
\draw[muted!22,very thin,xstep=1,ystep=1] (-4,-3) grid (5,4);
\draw[-latex,muted] (-4,0)--(5.35,0) node[right,font=\scriptsize]{$x$};
\draw[-latex,muted] (0,-3)--(0,4.35) node[above,font=\scriptsize]{$y$};
\node[font=\tiny,muted,below] at (-4,0) {-4};
\node[font=\tiny,muted,below] at (-2,0) {-2};
\node[font=\tiny,muted,below] at (2,0) {2};
\node[font=\tiny,muted,below] at (4,0) {4};
\node[font=\tiny,muted,left] at (0,-2) {-2};
\node[font=\tiny,muted,left] at (0,2) {2};
\begin{scope}\clip (-4,-3) rectangle (5,4);
\draw[accent,line width=1.1pt] (-3.0,-1.0)--(1.0,3.0);
\draw[accent,line width=1.1pt] (1.0,0.0)--(4.0,-2.0);
\end{scope}
\filldraw[fill=white,draw=accent,line width=0.9pt] (-3.0,-1.0) circle (2.2pt);
\fill[accent] (1.0,3.0) circle (2.2pt);
\filldraw[fill=white,draw=accent,line width=0.9pt] (1.0,0.0) circle (2.2pt);
\filldraw[fill=white,draw=accent,line width=0.9pt] (4.0,-2.0) circle (2.2pt);
\end{tikzpicture}
  \task \begin{tikzpicture}[x=0.27cm,y=0.27cm,baseline=(current bounding box.center)]
\draw[muted!22,very thin,xstep=1,ystep=1] (-4,-2) grid (4,3);
\draw[-latex,muted] (-4,0)--(4.35,0) node[right,font=\scriptsize]{$x$};
\draw[-latex,muted] (0,-2)--(0,3.35) node[above,font=\scriptsize]{$y$};
\node[font=\tiny,muted,below] at (-4,0) {-4};
\node[font=\tiny,muted,below] at (-2,0) {-2};
\node[font=\tiny,muted,below] at (2,0) {2};
\node[font=\tiny,muted,left] at (0,-2) {-2};
\node[font=\tiny,muted,left] at (0,2) {2};
\begin{scope}\clip (-4,-2) rectangle (4,3);
\draw[accent,line width=1.1pt] (-2.0,2.0)--(0.0,-1.0)--(2.0,2.0);
\draw[accent,line width=1.1pt,-latex] (-2.0,2.0)--(-2.467,2.700);
\draw[accent,line width=1.1pt,-latex] (2.0,2.0)--(2.467,2.700);
\end{scope}
\end{tikzpicture}
\end{tasks}\recorda{Domini: valors de $x$ que tenen imatge (projecció sobre l'eix $X$). Recorregut: valors que pren la $y$ (projecció sobre l'eix $Y$). Punt ple: $[\ ]$; punt buit: $(\ )$; fletxa: $\pm\infty$.}

\ex{1} Calcula el domini de:\begin{tasks}(4)
  \task $f(x)=3x^2-2x+1$
  \task $g(x)=\dfrac{4}{x-1}$
  \task $h(x)=\sqrt{x-2}$
  \task $k(x)=\dfrac{x+3}{x+5}$
\end{tasks}

\ex{2} Escriu el domini i el recorregut d'aquestes funcions:\begin{tasks}(3)
  \task \begin{tikzpicture}[x=0.25cm,y=0.25cm,baseline=(current bounding box.center)]
\draw[muted!22,very thin,xstep=1,ystep=1] (-6,-3) grid (6,5);
\draw[-latex,muted] (-6,0)--(6.35,0) node[right,font=\scriptsize]{$x$};
\draw[-latex,muted] (0,-3)--(0,5.35) node[above,font=\scriptsize]{$y$};
\node[font=\tiny,muted,below] at (-6,0) {-6};
\node[font=\tiny,muted,below] at (-4,0) {-4};
\node[font=\tiny,muted,below] at (-2,0) {-2};
\node[font=\tiny,muted,below] at (2,0) {2};
\node[font=\tiny,muted,below] at (4,0) {4};
\node[font=\tiny,muted,left] at (0,-2) {-2};
\node[font=\tiny,muted,left] at (0,2) {2};
\node[font=\tiny,muted,left] at (0,4) {4};
\begin{scope}\clip (-6,-3) rectangle (6,5);
\draw[accent,line width=1.1pt] (-5.0,1.0)--(-2.0,4.0);
\draw[accent,line width=1.1pt] (1.0,-2.0)--(4.0,0.0);
\draw[accent,line width=1.1pt,-latex] (4.0,0.0)--(5.800,1.200);
\end{scope}
\fill[accent] (-5.0,1.0) circle (2.2pt);
\fill[accent] (-2.0,4.0) circle (2.2pt);
\fill[accent] (1.0,-2.0) circle (2.2pt);
\end{tikzpicture}
  \task \begin{tikzpicture}[x=0.27cm,y=0.27cm,baseline=(current bounding box.center)]
\draw[muted!22,very thin,xstep=1,ystep=1] (-5,-4) grid (6,3);
\draw[-latex,muted] (-5,0)--(6.35,0) node[right,font=\scriptsize]{$x$};
\draw[-latex,muted] (0,-4)--(0,3.35) node[above,font=\scriptsize]{$y$};
\node[font=\tiny,muted,below] at (-4,0) {-4};
\node[font=\tiny,muted,below] at (-2,0) {-2};
\node[font=\tiny,muted,below] at (2,0) {2};
\node[font=\tiny,muted,below] at (4,0) {4};
\node[font=\tiny,muted,left] at (0,-4) {-4};
\node[font=\tiny,muted,left] at (0,-2) {-2};
\node[font=\tiny,muted,left] at (0,2) {2};
\begin{scope}\clip (-5,-4) rectangle (6,3);
\draw[accent,line width=1.1pt] (-4.0,-3.0)--(-1.0,0.0);
\draw[accent,line width=1.1pt] (-1.0,2.0)--(3.0,2.0);
\end{scope}
\fill[accent] (-4.0,-3.0) circle (2.2pt);
\filldraw[fill=white,draw=accent,line width=0.9pt] (-1.0,0.0) circle (2.2pt);
\fill[accent] (-1.0,2.0) circle (2.2pt);
\filldraw[fill=white,draw=accent,line width=0.9pt] (3.0,2.0) circle (2.2pt);
\fill[accent] (5.0,-1.0) circle (2.2pt);
\end{tikzpicture}
  \task \begin{tikzpicture}[x=0.27cm,y=0.27cm,baseline=(current bounding box.center)]
\draw[muted!22,very thin,xstep=1,ystep=1] (-5,-3) grid (5,3);
\draw[-latex,muted] (-5,0)--(5.35,0) node[right,font=\scriptsize]{$x$};
\draw[-latex,muted] (0,-3)--(0,3.35) node[above,font=\scriptsize]{$y$};
\node[font=\tiny,muted,below] at (-4,0) {-4};
\node[font=\tiny,muted,below] at (-2,0) {-2};
\node[font=\tiny,muted,below] at (2,0) {2};
\node[font=\tiny,muted,below] at (4,0) {4};
\node[font=\tiny,muted,left] at (0,-2) {-2};
\node[font=\tiny,muted,left] at (0,2) {2};
\begin{scope}\clip (-5,-3) rectangle (5,3);
\draw[accent,line width=1.1pt] (-3.0,-1.0)--(0.0,2.0);
\draw[accent,line width=1.1pt,-latex] (-3.0,-1.0)--(-4.800,-2.800);
\draw[accent,line width=1.1pt] (0.0,-2.0)--(3.0,1.0);
\draw[accent,line width=1.1pt,-latex] (3.0,1.0)--(4.800,2.800);
\end{scope}
\filldraw[fill=white,draw=accent,line width=0.9pt] (0.0,2.0) circle (2.2pt);
\filldraw[fill=white,draw=accent,line width=0.9pt] (0.0,-2.0) circle (2.2pt);
\end{tikzpicture}
\end{tasks}

\ex{2} Calcula el domini de:\begin{tasks}(4)
  \task $f(x)=\dfrac{1}{x^2-4}$
  \task $g(x)=\sqrt{6-2x}$
  \task $h(x)=\dfrac{x}{x^2+1}$
  \task $k(x)=\dfrac{5}{3x-9}$
\end{tasks}

\ex{2} Quin domini tenen aquestes funcions en el seu context?\begin{tasks}(1)
  \task El preu de $x$ bolígrafs a 0,80 € la unitat.
  \task L'àrea d'un quadrat de costat $x$: $A(x)=x^2$.
  \task L'altura d'una pilota, $h(t)=20t-5t^2$, mentre està en l'aire ($t$ en segons).
\end{tasks}

\ex{3} Dibuixa la gràfica d'una funció que tinga domini $[-3,\,4]$ i recorregut $[-2,\,5]$.

\ex{3} Calcula el domini de $f(x)=\dfrac{\sqrt{x+1}}{x-3}$.


% =====================================================
\seccio{5}{Creixement i decreixement}

\ex{1}\label{ex:mono1} Indica els intervals on aquestes funcions són creixents, decreixents o constants:\begin{tasks}(3)
  \task \begin{tikzpicture}[x=0.27cm,y=0.27cm,baseline=(current bounding box.center)]
\draw[muted!22,very thin,xstep=1,ystep=1] (-6,-3) grid (5,4);
\draw[-latex,muted] (-6,0)--(5.35,0) node[right,font=\scriptsize]{$x$};
\draw[-latex,muted] (0,-3)--(0,4.35) node[above,font=\scriptsize]{$y$};
\node[font=\tiny,muted,below] at (-6,0) {-6};
\node[font=\tiny,muted,below] at (-4,0) {-4};
\node[font=\tiny,muted,below] at (-2,0) {-2};
\node[font=\tiny,muted,below] at (2,0) {2};
\node[font=\tiny,muted,below] at (4,0) {4};
\node[font=\tiny,muted,left] at (0,-2) {-2};
\node[font=\tiny,muted,left] at (0,2) {2};
\begin{scope}\clip (-6,-3) rectangle (5,4);
\draw[accent,line width=1.1pt] (-5.0,-2.0)--(-2.0,3.0)--(1.0,-1.0)--(4.0,2.0);
\end{scope}
\fill[accent] (-5.0,-2.0) circle (2.2pt);
\fill[accent] (4.0,2.0) circle (2.2pt);
\end{tikzpicture}
  \task \begin{tikzpicture}[x=0.27cm,y=0.27cm,baseline=(current bounding box.center)]
\draw[muted!22,very thin,xstep=1,ystep=1] (-5,-4) grid (5,5);
\draw[-latex,muted] (-5,0)--(5.35,0) node[right,font=\scriptsize]{$x$};
\draw[-latex,muted] (0,-4)--(0,5.35) node[above,font=\scriptsize]{$y$};
\node[font=\tiny,muted,below] at (-4,0) {-4};
\node[font=\tiny,muted,below] at (-2,0) {-2};
\node[font=\tiny,muted,below] at (2,0) {2};
\node[font=\tiny,muted,below] at (4,0) {4};
\node[font=\tiny,muted,left] at (0,-4) {-4};
\node[font=\tiny,muted,left] at (0,-2) {-2};
\node[font=\tiny,muted,left] at (0,2) {2};
\node[font=\tiny,muted,left] at (0,4) {4};
\begin{scope}\clip (-5,-4) rectangle (5,5);
\draw[accent,line width=1.1pt] (-4.0,4.0)--(-1.0,1.0)--(2.0,1.0)--(4.0,-3.0);
\end{scope}
\fill[accent] (-4.0,4.0) circle (2.2pt);
\fill[accent] (4.0,-3.0) circle (2.2pt);
\end{tikzpicture}
  \task \begin{tikzpicture}[x=0.27cm,y=0.27cm,baseline=(current bounding box.center)]
\draw[muted!22,very thin,xstep=1,ystep=1] (-3,-4) grid (5,2);
\draw[-latex,muted] (-3,0)--(5.35,0) node[right,font=\scriptsize]{$x$};
\draw[-latex,muted] (0,-4)--(0,2.35) node[above,font=\scriptsize]{$y$};
\node[font=\tiny,muted,below] at (-2,0) {-2};
\node[font=\tiny,muted,below] at (2,0) {2};
\node[font=\tiny,muted,below] at (4,0) {4};
\node[font=\tiny,muted,left] at (0,-4) {-4};
\node[font=\tiny,muted,left] at (0,-2) {-2};
\begin{scope}\clip (-3,-4) rectangle (5,2);
\draw[accent,line width=1.1pt] (-1.0,-3.0)--(1.0,1.0)--(3.0,-1.0);
\draw[accent,line width=1.1pt,-latex] (-1.0,-3.0)--(-1.300,-3.600);
\draw[accent,line width=1.1pt,-latex] (3.0,-1.0)--(4.800,-2.800);
\end{scope}
\end{tikzpicture}
\end{tasks}\recorda{Llegim la gràfica d'esquerra a dreta. Els intervals es donen amb valors de la $x$ i són oberts.}

\ex{1} Sense dibuixar-les, digues si aquestes funcions són creixents, decreixents o constants:\begin{tasks}(4)
  \task $y=4x-1$
  \task $y=-3x+2$
  \task $y=5$
  \task $y=\tfrac{x}{2}$
\end{tasks}

\ex{2} Estudia el creixement i el decreixement d'aquestes funcions:\begin{tasks}(2)
  \task \begin{tikzpicture}[x=0.3cm,y=0.3cm,baseline=(current bounding box.center)]
\draw[muted!22,very thin,xstep=1,ystep=1] (-5,-3) grid (4,3);
\draw[-latex,muted] (-5,0)--(4.35,0) node[right,font=\scriptsize]{$x$};
\draw[-latex,muted] (0,-3)--(0,3.35) node[above,font=\scriptsize]{$y$};
\node[font=\tiny,muted,below] at (-4,0) {-4};
\node[font=\tiny,muted,below] at (-2,0) {-2};
\node[font=\tiny,muted,below] at (2,0) {2};
\node[font=\tiny,muted,left] at (0,-2) {-2};
\node[font=\tiny,muted,left] at (0,2) {2};
\begin{scope}\clip (-5,-3) rectangle (4,3);
\draw[accent,line width=1.1pt] (-3.0,2.0)--(-1.0,-2.0)--(1.0,-2.0)--(3.0,2.0);
\draw[accent,line width=1.1pt,-latex] (-3.0,2.0)--(-3.300,2.600);
\end{scope}
\fill[accent] (3.0,2.0) circle (2.2pt);
\end{tikzpicture}
  \task \begin{tikzpicture}[x=0.3cm,y=0.3cm,baseline=(current bounding box.center)]
\draw[muted!22,very thin,xstep=1,ystep=1] (-5,-3) grid (6,3);
\draw[-latex,muted] (-5,0)--(6.35,0) node[right,font=\scriptsize]{$x$};
\draw[-latex,muted] (0,-3)--(0,3.35) node[above,font=\scriptsize]{$y$};
\node[font=\tiny,muted,below] at (-4,0) {-4};
\node[font=\tiny,muted,below] at (-2,0) {-2};
\node[font=\tiny,muted,below] at (2,0) {2};
\node[font=\tiny,muted,below] at (4,0) {4};
\node[font=\tiny,muted,left] at (0,-2) {-2};
\node[font=\tiny,muted,left] at (0,2) {2};
\begin{scope}\clip (-5,-3) rectangle (6,3);
\draw[accent,line width=1.1pt] (-4.0,-1.0)--(-1.0,2.0);
\draw[accent,line width=1.1pt] (-1.0,-2.0)--(2.0,1.0)--(5.0,-1.0);
\end{scope}
\fill[accent] (-4.0,-1.0) circle (2.2pt);
\filldraw[fill=white,draw=accent,line width=0.9pt] (-1.0,2.0) circle (2.2pt);
\fill[accent] (-1.0,-2.0) circle (2.2pt);
\fill[accent] (5.0,-1.0) circle (2.2pt);
\end{tikzpicture}
\end{tasks}

\ex{2} Amb la gràfica de la temperatura de l'exercici \ref{ex:temp}, escriu els intervals de temps en què la temperatura puja i en què baixa.

\ex{3} En Pau diu que la funció de l'apartat a) de l'exercici \ref{ex:mono1} «creix en $(-2,\,3)$». Quin error ha comés?

\ex{3} Dibuixa la gràfica d'una funció de domini $\mathbb{R}$ que siga decreixent en $(-\infty,\,-1)$, constant en $(-1,\,2)$ i creixent en $(2,\,+\infty)$.


% =====================================================
\seccio{6}{Màxims i mínims}

\ex{1} Troba els màxims i mínims relatius d'aquestes funcions:\begin{tasks}(2)
  \task \begin{tikzpicture}[x=0.3cm,y=0.3cm,baseline=(current bounding box.center)]
\draw[muted!22,very thin,xstep=1,ystep=1] (-6,-3) grid (6,5);
\draw[-latex,muted] (-6,0)--(6.35,0) node[right,font=\scriptsize]{$x$};
\draw[-latex,muted] (0,-3)--(0,5.35) node[above,font=\scriptsize]{$y$};
\node[font=\tiny,muted,below] at (-6,0) {-6};
\node[font=\tiny,muted,below] at (-4,0) {-4};
\node[font=\tiny,muted,below] at (-2,0) {-2};
\node[font=\tiny,muted,below] at (2,0) {2};
\node[font=\tiny,muted,below] at (4,0) {4};
\node[font=\tiny,muted,left] at (0,-2) {-2};
\node[font=\tiny,muted,left] at (0,2) {2};
\node[font=\tiny,muted,left] at (0,4) {4};
\begin{scope}\clip (-6,-3) rectangle (6,5);
\draw[accent,line width=1.1pt] (-5.0,1.0)--(-3.0,4.0)--(0.0,-2.0)--(3.0,3.0)--(5.0,0.0);
\end{scope}
\fill[accent] (-5.0,1.0) circle (2.2pt);
\fill[accent] (5.0,0.0) circle (2.2pt);
\end{tikzpicture}
  \task \begin{tikzpicture}[x=0.3cm,y=0.3cm,baseline=(current bounding box.center)]
\draw[muted!22,very thin,xstep=1,ystep=1] (-5,-4) grid (4,4);
\draw[-latex,muted] (-5,0)--(4.35,0) node[right,font=\scriptsize]{$x$};
\draw[-latex,muted] (0,-4)--(0,4.35) node[above,font=\scriptsize]{$y$};
\node[font=\tiny,muted,below] at (-4,0) {-4};
\node[font=\tiny,muted,below] at (-2,0) {-2};
\node[font=\tiny,muted,below] at (2,0) {2};
\node[font=\tiny,muted,left] at (0,-4) {-4};
\node[font=\tiny,muted,left] at (0,-2) {-2};
\node[font=\tiny,muted,left] at (0,2) {2};
\begin{scope}\clip (-5,-4) rectangle (4,4);
\draw[accent,line width=1.1pt] (-3.0,-1.0)--(-1.0,3.0)--(2.0,-3.0);
\draw[accent,line width=1.1pt,-latex] (-3.0,-1.0)--(-4.300,-3.600);
\draw[accent,line width=1.1pt,-latex] (2.0,-3.0)--(2.300,-3.600);
\end{scope}
\end{tikzpicture}
\end{tasks}\recorda{Màxim relatiu: la funció passa de creixent a decreixent. Mínim relatiu: passa de decreixent a creixent. S'escriuen com a punts $(x,\,y)$.}

\ex{2} Troba els màxims i mínims, relatius i absoluts, d'aquestes funcions:\begin{tasks}(2)
  \task \begin{tikzpicture}[x=0.3cm,y=0.3cm,baseline=(current bounding box.center)]
\draw[muted!22,very thin,xstep=1,ystep=1] (-5,-3) grid (6,4);
\draw[-latex,muted] (-5,0)--(6.35,0) node[right,font=\scriptsize]{$x$};
\draw[-latex,muted] (0,-3)--(0,4.35) node[above,font=\scriptsize]{$y$};
\node[font=\tiny,muted,below] at (-4,0) {-4};
\node[font=\tiny,muted,below] at (-2,0) {-2};
\node[font=\tiny,muted,below] at (2,0) {2};
\node[font=\tiny,muted,below] at (4,0) {4};
\node[font=\tiny,muted,left] at (0,-2) {-2};
\node[font=\tiny,muted,left] at (0,2) {2};
\begin{scope}\clip (-5,-3) rectangle (6,4);
\draw[accent,line width=1.1pt] (-4.0,-2.0)--(-1.0,3.0)--(1.0,1.0)--(3.0,3.0)--(5.0,-1.0);
\end{scope}
\fill[accent] (-4.0,-2.0) circle (2.2pt);
\filldraw[fill=white,draw=accent,line width=0.9pt] (5.0,-1.0) circle (2.2pt);
\end{tikzpicture}
  \task \begin{tikzpicture}[x=0.3cm,y=0.3cm,baseline=(current bounding box.center)]
\draw[muted!22,very thin,xstep=1,ystep=1] (-4,-2) grid (4,5);
\draw[-latex,muted] (-4,0)--(4.35,0) node[right,font=\scriptsize]{$x$};
\draw[-latex,muted] (0,-2)--(0,5.35) node[above,font=\scriptsize]{$y$};
\node[font=\tiny,muted,below] at (-4,0) {-4};
\node[font=\tiny,muted,below] at (-2,0) {-2};
\node[font=\tiny,muted,below] at (2,0) {2};
\node[font=\tiny,muted,left] at (0,-2) {-2};
\node[font=\tiny,muted,left] at (0,2) {2};
\node[font=\tiny,muted,left] at (0,4) {4};
\begin{scope}\clip (-4,-2) rectangle (4,5);
\draw[accent,line width=1.1pt] (-3.0,4.0)--(0.0,-1.0)--(3.0,4.0);
\end{scope}
\filldraw[fill=white,draw=accent,line width=0.9pt] (-3.0,4.0) circle (2.2pt);
\filldraw[fill=white,draw=accent,line width=0.9pt] (3.0,4.0) circle (2.2pt);
\end{tikzpicture}
\end{tasks}

\ex{2} Fes una taula de valors de $f(x)=x^2-6x+5$ per a $x=0,\,1,\ldots,\,6$ i representa-la. On té el mínim? Quins són els punts de tall amb els eixos?

\ex{2} L'altura en metres d'un coet d'aigua és $h(t)=40t-5t^2$, amb $t$ en segons. Fes una taula per a $t=0,\,2,\,4,\,6,\,8$. Quina altura màxima assoleix i quan? Quan torna a terra?

\ex{3} Dibuixa una funció de domini $\mathbb{R}$ amb un màxim relatiu en $(-2,\,3)$ i un mínim relatiu en $(1,\,-1)$, que no tinga ni màxim ni mínim absoluts.

\ex{3} Raona si són vertaderes o falses:\begin{tasks}(1)
  \task Un màxim absolut sempre és també un màxim relatiu.
  \task Si una funció contínua té dos màxims relatius, entre ells hi ha d'haver almenys un mínim relatiu.
  \task Si una funció és creixent en tot el domini, no té màxims ni mínims relatius.
\end{tasks}


% =====================================================
\seccio{7}{Continuïtat}

\ex{1} Digues si aquestes funcions són contínues. Si no ho són, indica on són discontínues i per què:\begin{tasks}(3)
  \task \begin{tikzpicture}[x=0.27cm,y=0.27cm,baseline=(current bounding box.center)]
\draw[muted!22,very thin,xstep=1,ystep=1] (-5,-3) grid (5,3);
\draw[-latex,muted] (-5,0)--(5.35,0) node[right,font=\scriptsize]{$x$};
\draw[-latex,muted] (0,-3)--(0,3.35) node[above,font=\scriptsize]{$y$};
\node[font=\tiny,muted,below] at (-4,0) {-4};
\node[font=\tiny,muted,below] at (-2,0) {-2};
\node[font=\tiny,muted,below] at (2,0) {2};
\node[font=\tiny,muted,below] at (4,0) {4};
\node[font=\tiny,muted,left] at (0,-2) {-2};
\node[font=\tiny,muted,left] at (0,2) {2};
\begin{scope}\clip (-5,-3) rectangle (5,3);
\draw[accent,line width=1.1pt] (-4.0,-2.0)--(-1.0,1.0)--(2.0,-1.0)--(4.0,2.0);
\end{scope}
\fill[accent] (-4.0,-2.0) circle (2.2pt);
\fill[accent] (4.0,2.0) circle (2.2pt);
\end{tikzpicture}
  \task \begin{tikzpicture}[x=0.27cm,y=0.27cm,baseline=(current bounding box.center)]
\draw[muted!22,very thin,xstep=1,ystep=1] (-5,-2) grid (5,4);
\draw[-latex,muted] (-5,0)--(5.35,0) node[right,font=\scriptsize]{$x$};
\draw[-latex,muted] (0,-2)--(0,4.35) node[above,font=\scriptsize]{$y$};
\node[font=\tiny,muted,below] at (-4,0) {-4};
\node[font=\tiny,muted,below] at (-2,0) {-2};
\node[font=\tiny,muted,below] at (2,0) {2};
\node[font=\tiny,muted,below] at (4,0) {4};
\node[font=\tiny,muted,left] at (0,-2) {-2};
\node[font=\tiny,muted,left] at (0,2) {2};
\begin{scope}\clip (-5,-2) rectangle (5,4);
\draw[accent,line width=1.1pt] (-4.0,2.0)--(0.0,2.0);
\draw[accent,line width=1.1pt] (0.0,-1.0)--(4.0,3.0);
\end{scope}
\fill[accent] (-4.0,2.0) circle (2.2pt);
\filldraw[fill=white,draw=accent,line width=0.9pt] (0.0,2.0) circle (2.2pt);
\fill[accent] (0.0,-1.0) circle (2.2pt);
\fill[accent] (4.0,3.0) circle (2.2pt);
\end{tikzpicture}
  \task \begin{tikzpicture}[x=0.27cm,y=0.27cm,baseline=(current bounding box.center)]
\draw[muted!22,very thin,xstep=1,ystep=1] (-5,-2) grid (5,5);
\draw[-latex,muted] (-5,0)--(5.35,0) node[right,font=\scriptsize]{$x$};
\draw[-latex,muted] (0,-2)--(0,5.35) node[above,font=\scriptsize]{$y$};
\node[font=\tiny,muted,below] at (-4,0) {-4};
\node[font=\tiny,muted,below] at (-2,0) {-2};
\node[font=\tiny,muted,below] at (2,0) {2};
\node[font=\tiny,muted,below] at (4,0) {4};
\node[font=\tiny,muted,left] at (0,-2) {-2};
\node[font=\tiny,muted,left] at (0,2) {2};
\node[font=\tiny,muted,left] at (0,4) {4};
\begin{scope}\clip (-5,-2) rectangle (5,5);
\draw[accent,line width=1.1pt] (-4.0,-1.0)--(1.0,4.0);
\draw[accent,line width=1.1pt] (1.0,4.0)--(4.0,1.0);
\end{scope}
\fill[accent] (-4.0,-1.0) circle (2.2pt);
\filldraw[fill=white,draw=accent,line width=0.9pt] (1.0,4.0) circle (2.2pt);
\filldraw[fill=white,draw=accent,line width=0.9pt] (1.0,4.0) circle (2.2pt);
\fill[accent] (4.0,1.0) circle (2.2pt);
\fill[accent] (1.0,1.0) circle (2.2pt);
\end{tikzpicture}
\end{tasks}\recorda{Una funció és contínua si es pot dibuixar sense alçar el llapis del paper.}

\ex{1} Un aparcament cobra 1,50 € per cada hora o fracció, fins a un màxim de 4 hores. Representa la funció preu i calcula quant es paga per 2 h 30 min. És contínua?

\ex{2} Donada $f(x)=\begin{cases}x+3 & \text{si } x<0\\ 3-x^2 & \text{si } x\ge0\end{cases}$, calcula $f(-2)$, $f(0)$ i $f(2)$. És contínua en $x=0$?

\ex{2} Donada $g(x)=\begin{cases}2x & \text{si } x\le1\\ x+3 & \text{si } x>1\end{cases}$, calcula $g(1)$ i digues si és contínua en $x=1$. Dibuixa-la.

\ex{2} Una empresa de missatgeria cobra 3 € per paquets de fins a 1 kg, 5 € per paquets de més d'1 kg fins a 3 kg i 8 € per paquets de més de 3 kg fins a 5 kg. Dibuixa la gràfica, escriu el domini i el recorregut i indica on és discontínua.

\ex{3} Troba el valor de $a$ perquè la funció $f(x)=\begin{cases}x^2+a & \text{si } x<2\\ 3x-1 & \text{si } x\ge2\end{cases}$ siga contínua.

\ex{3} On són discontínues aquestes funcions? Per què?\begin{tasks}(2)
  \task $f(x)=\dfrac{1}{x-2}$
  \task $g(x)=\dfrac{3}{x^2-9}$
\end{tasks}


% =====================================================
\seccio{8}{Problemes amb funcions}

\ex{1} Un taxi cobra 3,50 € per la baixada de bandera i 1,10 € per quilòmetre. Escriu la fórmula del preu, calcula quant costa un trajecte de 12 km i quants quilòmetres hem fet si hem pagat 25,50 €.\recorda{Identifica les variables (què és la $x$ i què és la $y$, amb les unitats), escriu la fórmula o dibuixa la gràfica i interpreta cada resultat en el context del problema.}

\ex{1} Una tarifa de mòbil A costa 10 € al mes més 0,05 € per minut, i una tarifa B costa 0,15 € per minut. Escriu les dues fórmules. Quina convé a qui parla 60 minuts al mes? I 150 minuts? Per a quants minuts costen el mateix?

\ex{2} Fes l'estudi complet d'aquesta funció: domini, recorregut, punts de tall, creixement, màxims i mínims i continuïtat.\par\smallskip\begin{center}\begin{tikzpicture}[x=0.4cm,y=0.4cm,baseline=(current bounding box.center)]
\draw[muted!22,very thin,xstep=1,ystep=1] (-6,-5) grid (6,4);
\draw[-latex,muted] (-6,0)--(6.35,0) node[right,font=\scriptsize]{$x$};
\draw[-latex,muted] (0,-5)--(0,4.35) node[above,font=\scriptsize]{$y$};
\node[font=\tiny,muted,below] at (-6,0) {-6};
\node[font=\tiny,muted,below] at (-4,0) {-4};
\node[font=\tiny,muted,below] at (-2,0) {-2};
\node[font=\tiny,muted,below] at (2,0) {2};
\node[font=\tiny,muted,below] at (4,0) {4};
\node[font=\tiny,muted,left] at (0,-4) {-4};
\node[font=\tiny,muted,left] at (0,-2) {-2};
\node[font=\tiny,muted,left] at (0,2) {2};
\begin{scope}\clip (-6,-5) rectangle (6,4);
\draw[accent,line width=1.1pt] (-5.0,0.0)--(-3.0,3.0)--(0.0,-2.0);
\draw[accent,line width=1.1pt] (0.0,-4.0)--(2.0,1.0)--(5.0,-1.0);
\end{scope}
\fill[accent] (-5.0,0.0) circle (2.2pt);
\fill[accent] (0.0,-2.0) circle (2.2pt);
\filldraw[fill=white,draw=accent,line width=0.9pt] (0.0,-4.0) circle (2.2pt);
\fill[accent] (5.0,-1.0) circle (2.2pt);
\end{tikzpicture}\end{center}

\ex{2} Un dipòsit té 20 L d'aigua. S'omple durant 5 minuts fins als 120 L, es manté així 10 minuts i després es buida del tot en 10 minuts. Dibuixa la gràfica volum--temps i estudia-la. A quina velocitat s'omple i a quina es buida?

\ex{2} Amb 24 m de tanca volem fer un corral rectangular. Escriu l'àrea en funció de la base $x$, fes una taula per a $x=1,\,3,\,6,\,9,\,11$ i digues quines dimensions donen l'àrea màxima.

\ex{2} Llancem una pilota cap amunt i la seua altura en metres és $h(t)=30t-5t^2$ ($t$ en segons). Fes una taula per a $t=0,\,1,\,2,\,3,\,4,\,6$, representa-la i respon: quina altura màxima assoleix? Quan torna a terra? Quin domini té sentit?

\ex{3} El gimnàs A cobra 30 € de quota mensual més 2 € per sessió; el gimnàs B no té quota però cobra 5 € per sessió. Escriu les dues funcions, representa-les en els mateixos eixos i digues a partir de quantes sessions al mes convé el gimnàs A.

\ex{3} La tarifa de l'aigua d'un poble és 0,50 €/m\textsuperscript{3} fins a 10 m\textsuperscript{3} i 0,80 €/m\textsuperscript{3} per cada m\textsuperscript{3} que passa de 10. Escriu la funció a trossos, calcula el preu de 15 m\textsuperscript{3} i digues si és contínua.

\ex{3} L'Anna ix de casa caminant a 5 km/h. Una hora després, el seu germà ix pel mateix camí amb bicicleta a 15 km/h. Escriu la distància a casa de cadascun en funció del temps (des que ix l'Anna) i troba quan i on l'atrapa.

\vfill
\begin{center}\color{muted}\small\faCheckCircle[regular]\ \ Ànim! Si et costa algun exercici, repassa els apunts interactius de la unitat.\end{center}

\end{document}
